\section{贝叶斯公式与条件分数分解}

解决逆问题的核心数学工具是\textbf{贝叶斯公式}，它将条件分数分解为边际分数和似然分数的和。

\subsection{条件分数的贝叶斯分解}

我们的目标是计算条件分数 $\nabla_{x_t} \log p(x_t | y)$，以便在去噪过程中引导生成。利用贝叶斯公式：

$$
p(x_t | y) = \frac{p(y | x_t) \cdot p(x_t)}{p(y)}
$$

两边取对数再求梯度：

$$
\nabla_{x_t} \log p(x_t | y) = \underbrace{\nabla_{x_t} \log p(x_t)}_{\text{边际分数（已知）}} + \underbrace{\nabla_{x_t} \log p(y | x_t)}_{\text{似然分数（待计算）}}
$$

\begin{importantbox}{条件分数 = 边际分数 + 似然分数}
\begin{itemize}
    \item \textbf{边际分数} $\nabla_{x_t} \log p(x_t)$：这正是预训练扩散模型已经学会预测的分数函数，可以直接从模型中获得。
    \item \textbf{似然分数} $\nabla_{x_t} \log p(y | x_t)$：衡量在当前中间状态 $x_t$ 下，观测 $y$ 出现的可能性有多大。这是需要额外计算（或近似）的部分。
\end{itemize}
因此，整个问题归结为：\textbf{如何在不训练额外网络的情况下，近似计算似然分数？}
\end{importantbox}

\subsection{与 Classifier Guidance 的对比}

\begin{knowledgebox}{统一视角下的已有方法}
之前学过的引导方法可以在此框架下统一理解：
\begin{itemize}
    \item \textbf{Classifier Guidance}：通过训练一个额外的分类器网络来直接估计似然分数 $\nabla_{x_t} \log p(y | x_t)$。
    \item \textbf{Classifier-Free Guidance}：用同一个网络联合学习条件和无条件噪声，隐式地通过两者的差异来计算似然分数。
\end{itemize}
这两种方法都需要训练。我们现在要做的是\textbf{无需训练地近似似然分数}。
\end{knowledgebox}

\subsection{本章小结}

贝叶斯分解告诉我们，条件生成只需要在预训练模型的边际分数之上，额外添加一个似然分数即可。剩下的核心挑战是：如何高效地近似这个似然分数。


